%---------- Inleiding ---------------------------------------------------------

\section{Inleiding}
\label{sec:inleiding}
Deze bachelorproef onderzoekt Master Data Management (MDM), multi-ERP-omgevingen en het concept van "golden records''. Specifiek richt het onderzoek zich op de casus van IPCOM, een bedrijf in de isolatiesector dat met meerdere ERP-systemen werkt. IPCOM heeft op dit moment het probleem dat verschillende entiteiten binnen de groep (zoals in Oostenrijk, Zwitserland en Noorwegen) dezelfde artikelen van dezelfde leveranciers kopen, maar dat deze door de verschillende bronsystemen niet eenvoudig te identificeren zijn als hetzelfde product. Hierdoor ontstaan inconsistenties in artikel- en leveranciersdata, wat een directe impact heeft op rapportering en besluitvorming.

Om dit probleem aan te pakken is IPCOM momenteel een pilot gestart met de commerciële MDM-tool Profisee. Dit roept een interessante nieuwe vraag op voor deze bachelorproef. In plaats van een theoretische oplossing te bedenken, wordt er een proof-of-concept (PoC) gebouwd in Microsoft Fabric die als objectieve benchmark dient. Hoe presteert een eigen, in-house gebouwde oplossing op basis van deterministische regels en machine learning in vergelijking met de Profisee-setup, wanneer exact dezelfde inputdata en mappings worden gebruikt? Hoe past zo'n MDM-tool in een datawarehouse-architectuur en hoe kan de datakwaliteit over de verschillende systemen heen worden gemeten?

De belangrijkste doelgroep wordt gevormd door data-engineers en business-gebruikers binnen de organisatie die werken met deze data in het datawarehouse en het rapportageplatform (bijvoorbeeld Power BI). Zij ondervinden dagelijks hinder van gefragmenteerde data en hebben nood aan een gecentraliseerde, betrouwbare bron van waarheid, namelijk een "golden record" per artikel en leverancier, die de datakwaliteit structureel verbetert.

De centrale onderzoeksvraag luidt als volgt: Hoe kan een in-house Master Data Management-oplossing in Microsoft Fabric worden ontworpen en geïmplementeerd voor artikel- en leveranciersdata, zodat deze kan dienen als objectieve benchmark tegenover een commerciële MDM-tool (Profisee) binnen een multi-ERP-omgeving?

De doelstelling van deze bachelorproef is eenduidig: een PoC bouwen voor een MDM-tool in Microsoft Fabric, specifiek gericht op de ERP-data van Oostenrijk, Zwitserland en Noorwegen. Op basis van de vergelijking tussen deze PoC en de Profisee-pilot kunnen concrete aanbevelingen worden geformuleerd over de toegevoegde waarde van commerciële pakketten versus in-house oplossingen voor data harmonisatie en matching.

\subsection{Deelvragen m.b.t probleemdomein}
\begin{itemize}
    \item Hoeveel duplicaten en conflicterende records komen voor over de geselecteerde ERP-systemen heen?
    \item Welke datakwaliteitsproblemen treden het meest op bij de geselecteerde artikel- en leveranciersdata (bijvoorbeeld het hoge percentage ontbrekende waarden)?
    \item Welke identificatievelden (zoals EAN en artikelnaam) zijn bruikbaar in de systemen om een correcte mapping op te zetten?
    \item Hoe vaak verandert een golden record de oorspronkelijke data bij welk specifiek bronsysteem?
\end{itemize}

\subsection{Deelvragen m.b.t oplossingsdomein}
\begin{itemize}
    \item Hoe presteert de in-house Fabric PoC ten opzichte van de commerciële Profisee-setup wanneer dezelfde mappings en rulesets worden gebruikt?
    \item Welke matching-technieken (deterministisch op basis van EAN versus fuzzy matching of AI op artikelnaam) leveren de hoogste nauwkeurigheid?
    \item Welke consolidatieregels (survivorship-logica) zorgen voor de meest consistente "golden records"?
    \item Hoe kan de publicatie van golden records worden ingericht zodat lokale id's correct verwijzen naar de golden records en alle entiteiten hiervan profiteren?
    \item Hoeveel verbetering in datakwaliteit wordt gerealiseerd en visueel meetbaar gemaakt na implementatie van de MDM-oplossing?
\end{itemize}





%---------- Stand van zaken ---------------------------------------------------

\section{Stand van zaken}%
\label{sec:stand-van-zaken}

\subsection{Master Data Management}
Master Data Management (MDM) wordt omschreven als een methode om master data te onderhouden, te integreren en te harmoniseren, met als doel consistente informatievoorziening over informatiesystemen heen te ondersteunen \autocite{Hikmawati}. Master data verwijst daarbij naar kerngegevens die organisatiebreed worden hergebruikt en die verschillende processen en afdelingen met elkaar verbinden. Vilminko-Heikkinen en Pekkola positioneren MDM hierbij als een organisatiebrede data-managementpraktijk die expliciet gericht is op het verbeteren van de kwaliteit van deze kerngegevens \autocite{VilminkoHeikkinenPekkola2019}.

De kernfunctie van MDM is het beheersen van master data zodat deze consistent, accuraat, actueel en relevant blijft over applicaties en organisatie-eenheden heen \autocite{Hikmawati}. Wanneer dit beheer tekortschiet, neemt het risico op onnauwkeurige of tegenstrijdige gegevens toe, wat de betrouwbaarheid van rapportering en besluitvorming ondermijnt \autocite{Hikmawati}. Zeker wanneer gegevens verspreid zijn over meerdere systemen, kunnen definities en formaten verschillen, waardoor fragmentatie ontstaat \autocite{VilminkoHeikkinenPekkola2019}. Dit maakt MDM naast een technisch ook een organisatorisch vraagstuk: meerdere teams beheren immers dezelfde kernentiteiten.

Vilminko-Heikkinen en Pekkola benaderen MDM niet als louter softwaretool, maar als een praktijk die aanpassingen in processen en werkwijzen vereist \autocite{VilminkoHeikkinenPekkola2019}. Het omvat beleid en infrastructuur om master data te capteren en te delen, met focus op consistentie en tijdigheid. Dit impliceert dat een MDM-initiatief onlosmakelijk verbonden is met \emph{data governance}: het vastleggen van afspraken en controlemechanismen \autocite{Hikmawati}. Uit onderzoek van Hikmawati blijkt dan ook dat MDM-initiatieven vaak leiden tot een hertekening van rollen om datakwaliteit structureel te borgen.

Governance uit zich concreet in rollen en verantwoordelijkheden. In een MDM-context zijn rollen zoals \emph{data owner} (eindverantwoordelijkheid) en \emph{data steward} (operationele bewaking) essentieel om eigenaarschap te organiseren \autocite{Hikmawati}. Deze verdeling maakt het mogelijk kwaliteitsproblemen toe te wijzen en op te volgen \autocite{VilminkoHeikkinenPekkola2019}. Zonder expliciet eigenaarschap blijft het risico bestaan dat master data opnieuw divergeert, zelfs bij aanwezigheid van een technisch MDM-platform \autocite{VilminkoHeikkinenPekkola2019}.

Voor IPCOM is dit relevant omdat product- en leveranciersgegevens uit meerdere ERP-systemen samenkomen. MDM kan bijdragen aan harmonisatie door definities en kwaliteitsregels vast te leggen \autocite{Hikmawati}. Echter, zoals de literatuur benadrukt, hangt succes af van organisatorische verankering. Binnen deze bachelorproef wordt MDM daarom onderzocht als methode om duplicatie te reduceren, mits de implementatie zowel technisch als organisatorisch wordt gedragen \autocite{Hikmawati}.

\subsection{Golden records en dataconsolidatie}
De multi-ERP-context van IPCOM impliceert dat master data verspreid is, waardoor informatie over één entiteit gefragmenteerd voorkomt \autocite{Elahi}. Elahi koppelt dit aan een hogere werkdruk, omdat gebruikers meerdere bronnen moeten consolideren. Een manuele aanpak vergroot bovendien het risico op fouten; een risico dat toeneemt naarmate het aantal bronsystemen groeit \autocite{Elahi}.

Een centraal concept om dit op te lossen is het \emph{golden record}: de enkelvoudige weergave die ontstaat nadat duplicaten zijn samengebracht en relevante informatie is behouden \autocite{Elahi}. Bij IPCOM kan dit gaan om een product waarvan attributen (zoals leveranciersnummer, prijs, of omschrijving) in verschillende systemen bestaan. Het golden record brengt per attribuut de geselecteerde waarde samen tot één consistente representatie voor analyse en rapportering \autocite{Elahi}.



Elahi onderscheidt drie stappen die hieraan voorafgaan. \emph{Data matching} groepeert records op basis van waarschijnlijkheid. Vervolgens kan \emph{deduplicatie} of \emph{data merging} de datasets combineren. De brug naar het golden record wordt gevormd door \emph{data survivorship}: het selecteren van de beste attributen uit duplicaten om één uniforme record te vormen \autocite{Elahi}.

Er zijn meerdere benaderingen voor survivorship mogelijk, afhankelijk van datakwaliteit \autocite{Elahi}:
\begin{itemize}
    \item \textbf{Manuele selectie:} Het meest volledige record kiezen (moeilijk schaalbaar).
    \item \textbf{Bronprioriteit:} Selectie op basis van de betrouwbaarheid van het bronsysteem.
    \item \textbf{Attribuutlogica:} Per veld de beste waarde kiezen (bijv. meest recente update).
    \item \textbf{Samengesteld (composite):} Waarden conditioneel overschrijven om informatie uit meerdere bronnen te combineren.
\end{itemize}

De operationalisering gebeurt via \emph{survivorship rules} (bijv. 'langste waarde wint' of 'laatste update wint'). Elahi benadrukt dat deze regels iteratief moeten worden bijgestuurd op basis van de uitkomst \autocite{Elahi}. In dit onderzoek wordt geëvalueerd welke regels passen bij de datakwaliteit van IPCOM: volstaat bronprioriteit, of is een complexe, attribuutgebaseerde benadering noodzakelijk?

\subsection{Multi-ERP-omgevingen en de rol van Master Data Management}
In organisaties met een \emph{multi-ERP-landschap} zijn master data en processen verdeeld. Maedche (2010) beschrijft dat integratie in grote ondernemingen lastig blijft wanneer het aantal en de diversiteit van ERP-systemen toeneemt. Hierdoor ontstaat een masterdata-probleem: dezelfde entiteiten bestaan in meerdere systemen met afwijkende definities \autocite{Maedche2010}.

MDM verschilt hierin fundamenteel van ERP. Waar ERP procesgericht is, beoogt MDM een \emph{enterprise-wide single version of the truth} die onafhankelijk is van specifieke applicaties \autocite{Maedche2010}. Omdat ERP en MDM elkaar functioneel kunnen overlappen, verschuift het ontwerpprobleem naar de complementariteit tussen beide in een complex landschap.

De uitdagingen zijn vooral zichtbaar in \emph{operationele} MDM-scenario’s. Standalone MDM-oplossingen kunnen integratieproblemen veroorzaken rond distributie en semantische mapping \autocite{Maedche2010}. Dit wijst erop dat multi-ERP-omgevingen niet alleen consolidatie vereisen, maar ook een integratieaanpak die synchronisatie beheersbaar maakt.

Maedche onderscheidt in deze context twee scenario's:
\begin{enumerate}
    \item \emph{Collaboration in corporate groups:} Een stertopologie met centraal beheer en lokale restricties.
    \item \emph{Heterogeneous system landscapes:} Een mesh-topologie zonder duidelijke 'master', waarbij consolidatie vooral dient voor rapportering \autocite{Maedche2010}.
\end{enumerate}

Om MDM te laten aansluiten stelt Maedche een ERP-centrische benadering voor via \emph{shared MDM capabilities}. Dit omvat een MDM-hub, integratiemiddleware en een ingebedde MDM-component in de ERP-systemen. Het doel is dat wijzigingen consistent en semantisch correct worden doorgegeven \autocite{Maedche2010}.

Voor IPCOM betekent dit dat de keuze voor een integratiepatroon cruciaal is. Afhankelijk van de behoefte aan lokale autonomie kan een hybride benadering worden onderzocht, waarbij centrale standaarden worden gecombineerd met lokale verrijking \autocite{Maedche2010}.

\subsection{Data quality en data governance binnen MDM}
Datakwaliteit is zowel een drijfveer als een evaluatiecriterium voor MDM. Kwaliteit wordt hierbij gezien als geschiktheid voor gebruik, gemeten via dimensies zoals accuraatheid, volledigheid, consistentie en tijdigheid \autocite{Hikmawati}. Problemen zoals duplicaten of ontbrekende waarden beïnvloeden direct de besluitvorming \autocite{Hikmawati}.

MDM beoogt deze problemen te reduceren, maar Vilminko-Heikkinen en Pekkola benadrukken dat datakwaliteit geen louter technisch resultaat is; het vereist een combinatie van processen en technologie \autocite{VilminkoHeikkinenPekkola2019}. Dit veronderstelt duidelijke afspraken over datacreatie en validatie, verankerd in \emph{data governance}.

Zonder heldere governance (rollen, eigenaarschap) wordt kwaliteitsborging bemoeilijkt \autocite{VilminkoHeikkinenPekkola2019}. Hikmawati bevestigt dat onduidelijk ownership en gebrek aan procedures de effectiviteit van MDM ondermijnen. Governance concretiseert wie beslissingsrecht heeft (data owner) en wie de kwaliteit bewaakt (data steward), vaak ondersteund door een \emph{data governance council} \autocite{Hikmawati}.

Het proces start doorgaans met \emph{data profiling} om problemen zichtbaar te maken, gevolgd door opschoning (cleaning) en synchronisatie. Deze stappen zijn pas effectief wanneer ze gekoppeld zijn aan meetbare kwaliteitsdoelen, zoals indicatoren voor de duplicatiegraad \autocite{Hikmawati}.

Voor IPCOM betekent dit dat de technische oplossing gepaard moet gaan met governance die lokale verantwoordelijkheid en centrale uniformiteit combineert. Door data ownership toe te wijzen, wordt datakwaliteit een bestuurbaar resultaat in plaats van een toevallige uitkomst \autocite{VilminkoHeikkinenPekkola2019}.

\subsection{Bestaande MDM-oplossingen en implementatiebenaderingen}
MDM-oplossingen variëren in architectuur, synchronisatiepatronen en governance-ondersteuning \autocite{Rengasamy2025}. Vooral in multi-ERP-omgevingen is de keuze tussen centralisatie en federatie bepalend.



\begin{itemize}
    \item \textbf{Gecentraliseerd:} Eén MDM-hub fungeert als de primaire bron.
    \item \textbf{Federated:} Eigenaarschap blijft bij businessfuncties, terwijl consistentie via afspraken wordt bewaakt. Dit is geschikt wanneer lokale autonomie zwaar weegt \autocite{Rengasamy2025}.
\end{itemize}

Daarnaast onderscheidt de literatuur drie MDM-patronen \autocite{Rengasamy2025}:
\begin{enumerate}
    \item \emph{Registry:} Focus op referenties en identifiers zonder dataverplaatsing.
    \item \emph{Coexistence:} Lokale opslag met synchronisatie naar een centrale standaard.
    \item \emph{Transaction Hub:} Real-time verwerking van updates als centrale autoriteit.
\end{enumerate}

Voor multi-ERP-contexten worden vaak Coexistence of Transaction Hubs aangeraden, omdat bronsystemen operationeel moeten blijven \autocite{Maedche2010}. Rengasamy (2025) adviseert een gefaseerde implementatie: starten met een afgebakend domein en iteratief uitbreiden. Dit sluit aan bij de visie dat datakwaliteit een continu proces is \autocite{Hikmawati}.

Voor IPCOM wordt een federated of coexistence-achtige benadering onderzocht. Dit maakt het mogelijk om centrale definities te combineren met operationele continuïteit in de ERP-systemen \autocite{Rengasamy2025}.

\subsection{Open uitdagingen en positionering van dit onderzoek}
MDM-initiatieven kennen zowel technische als organisatorische uitdagingen \autocite{Pansara2021}.
Ten eerste is er \textbf{data governance}: zonder duidelijke rollen en procedures kunnen inconsistenties opnieuw in systemen sluipen \autocite{Pansara2021}. De vraag is hoe governance consistentie kan afdwingen zonder lokale werking te blokkeren \autocite{VilminkoHeikkinenPekkola2019}.

Ten tweede vormt \textbf{data-integratie} een risico. Synchronisatieverschillen en semantische mismatches tussen MDM en ERP kunnen leiden tot fouten \autocite{Maedche2010}. Ook bij het samenvoegen van records (survivorship) kan context verloren gaan \autocite{Pansara2021}.

Ten derde vereisen \textbf{data-standaarden} afstemming tussen afdelingen, wat in gedecentraliseerde contexten moeilijk is.

Deze bachelorproef positioneert zich als toegepast onderzoek. MDM-concepten, matching-technieken en survivorship-strategieën worden vertaald naar een proof-of-concept voor de multi-ERP-casus van IPCOM. Het doel is om deze in-house oplossing in te zetten als een objectieve benchmark tegenover een commerciële MDM-pilot (Profisee). Hierbij wordt geëvalueerd hoe duplicatie effectief kan worden gereduceerd via deterministische en AI-gebaseerde matching, en of een in-house setup kwalitatief kan concurreren met een kant-en-klaar platform.


%---------- Methodologie ------------------------------------------------------
\section{Methodologie}%
\label{sec:methodologie}

\subsection{Onderzoeksopzet}
Deze bachelorproef is \textbf{toegepast onderzoek} binnen de bedrijfscontext van IPCOM en wordt uitgevoerd als \textbf{case study} met een \textbf{PoC} als technisch eindresultaat. De centrale onderzoeksvraag uit Sectie~\ref{sec:inleiding} wordt beantwoord door (1) literatuurstudie, (2) requirements- en contextanalyse via directe afstemming met de interne data engineer, (3) ontwerp en implementatie van een PoC in Microsoft Fabric, waarbij de scope bewust beperkt wordt tot artikel- en leveranciersdata uit drie specifieke ERP-systemen (Oostenrijk, Zwitserland en Noorwegen). Tot slot volgt (4) een evaluatie waarbij deze in-house PoC fungeert als objectieve benchmark tegenover een lopende pilot met het commerciële MDM-platform Profisee.

\subsection{Onderzoekstechnieken}
De gebruikte onderzoekstechnieken zijn:
\begin{itemize}
  \item \textbf{Literatuurstudie}: concepten, implementatiebenaderingen en open uitdagingen rond MDM, multi-ERP, golden records en datakwaliteit.
  \item \textbf{Expert-consultatie en documentanalyse}: directe afstemming met de co-promotor (data engineer) als hoofd-stakeholder en het hergebruiken van reeds intern opgestelde rule-sets en datamappings.
  \item \textbf{Data-analyse en data profiling}: bepalen van een objectieve baseline (zoals het huidige percentage ontbrekende velden) op representatieve datasets uit de geselecteerde bronsystemen.
  \item \textbf{Ontwerp en implementatie PoC}: data harmonisatie, deterministische en AI-gestuurde matching, survivorship rules en golden record-opbouw in Microsoft Fabric.
  \item \textbf{Evaluatie (Benchmark)}: kwantitatieve en kwalitatieve vergelijking van de Fabric PoC-output met zowel de originele baseline als de resultaten van de Profisee-oplossing.
\end{itemize}

\subsection{Dataverzameling en requirements}
 Er is directe toegang verleend tot de benodigde databases. Dit versnelt de dataverzameling aanzienlijk. De scope is afgebakend tot entiteiten van het type \emph{Article} en \emph{Supplier}, aangezien deze de grootste kans op overlap vertonen over de verschillende entiteiten heen. Klantendata wordt buiten beschouwing gelaten vanwege het sterk lokale karakter. 

Er worden geen brede stakeholder-interviews gehouden; de requirements, pijnpunten en reeds bestaande mappings (bijvoorbeeld vertalingen van landcodes en lokale veldnamen) worden rechtstreeks gecapteerd via de interne data engineer. Dit vormt de fundering voor het canonical model.

\subsection{Technische uitvoering van de PoC}
De PoC implementeert een reproduceerbare pipeline in Microsoft Fabric met de volgende stappen:
\begin{enumerate}
  \item \textbf{Landing}: inladen van datasnapshots uit de ERP-systemen van Zwitserland, Oostenrijk en Noorwegen.
  \item \textbf{Harmonisatie}: toepassen van de bestaande mappings om data te vertalen naar één uniform canonical model.
  \item \textbf{Matching}: groeperen van records. Dit verloopt in twee fasen om de benchmark te voeden:
  \begin{itemize}
      \item \emph{Deterministisch}: harde regels, zoals een match op basis van identieke EAN-codes in combinatie met hetzelfde leveranciers-ID.
      \item \emph{Fuzzy/AI-gebaseerd}: experimenteel matchen op artikelnamen als aanvulling, om te evalueren of dit superieure resultaten oplevert ten opzichte van de standaard logica.
  \end{itemize}
  \item \textbf{Survivorship}: toepassen van regels om het ultieme golden record te genereren, waarbij ook lokale referentie-ID's worden bijgehouden zodat bronsystemen via "foreign keys" kunnen terugverwijzen.
  \item \textbf{Golden record store \& Publicatie}: opslaan van de resultaten met uitgebreide auditvelden voor downstream consumptie in Power BI.
\end{enumerate}

\subsection{Evaluatiecriteria}
De evaluatie gebeurt door een vergelijking op basis van meetbare indicatoren:
\begin{itemize}
  \item \textbf{Benchmark Profisee}: presteert de opgezette logica in Fabric vergelijkbaar, beter of slechter dan de commerciële opzet wanneer dezelfde inputdata wordt gebruikt?
  \item \textbf{Impact op brondata}: hoe vaak en in welke mate wijzigt het gegenereerde golden record de originele data in een specifiek bronsysteem?
  \item \textbf{Volledigheid en Consistentie}: afname van het aantal ontbrekende waarden en conflicten tussen systemen voor dezelfde entiteit.
  \item \textbf{Traceerbaarheid}: transparantie in welke bronrecords en rules hebben geleid tot een specifieke golden record-waarde.
\end{itemize}

\subsection{Tools en technologie (initiële keuze)}
De architectuur leunt op het Microsoft-ecosysteem, met \textbf{Microsoft Fabric} (Lakehouse, Dataflows Gen2, Notebooks in Python/SQL) als integratieomgeving. \textbf{Power BI} wordt ingezet voor visuele validatie en het presenteren van de datakwaliteit en benchmark-metrieken.

\subsection{Planning (14 weken, wekelijkse cyclus)}
Door de interne toegang tot data en bestaande mappings kan de technische opstart versneld worden uitgevoerd.

\begin{table*}[t]
\centering
\small
\setlength{\tabcolsep}{4pt}
\renewcommand{\arraystretch}{1.15}
\caption{Planning en deliverables (14 vrijdagen)}
\label{tab:planning}
\begin{tabularx}{\textwidth}{
  >{\raggedright\arraybackslash}p{0.07\textwidth}
  >{\raggedright\arraybackslash}X
  >{\raggedright\arraybackslash}X
}
\toprule
\textbf{Week} & \textbf{Activiteit} & \textbf{Deliverable} \\
\midrule
1  & Kick-off: scope-afbakening (Profisee benchmark) \& ERP-selectie (CH, AT, NO)
   & Scope + afspraken gedocumenteerd \\
2  & Data-extractie (Article/Supplier) + verzamelen bestaande mappings
   & Dataset snapshot + mappingdocument \\
3  & Data landing in Fabric + eerste harmonisatie o.b.v. mapping
   & Werkende landing \& harmonisatielaag \\
4  & Data profiling + meten van baseline metrieken (bijv. $\%$ lege velden)
   & Baseline DQ-rapport \\
5  & Ontwerp PoC-architectuur \& voorbereiding rule-engine
   & Architectuurdiagram \\
6  & Implementatie matching fase 1: deterministische rules (EAN/Supplier)
   & Werkende hard-rule set \\
7  & Implementatie matching fase 2: Fuzzy matching / AI op Artikelnaam
   & AI-matching testresultaten \\
8  & Implementatie survivorship rules \& opbouw golden records
   & Golden record generator \\
9  & Traceerbaarheid inbouwen (logging van gekozen bronwaarden)
   & Traceability-output + auditvelden \\
10 & Evaluatie: analyse "Impact op brondata" (wijzigingsfrequentie)
   & Data-impact analyse \\
11 & Vergelijking en benchmark resultaten tegenover de Profisee-oplossing
   & Benchmark-rapportage \\
12 & Publicatie consumptielaag + Power BI validatie dashboards
   & Power BI proof + visualisaties \\
13 & Stakeholder review met co-promotor (inzichten en werkbaarheid)
   & Reviewverslag \\
14 & Finalisatie scriptie + formuleren advies (Fabric vs. Profisee)
   & Eindrapport + aanbevelingen \\
\bottomrule
\end{tabularx}
\end{table*}


%---------- Verwachte resultaten ----------------------------------------------
\section{Verwacht resultaat, conclusie}%
\label{sec:verwachte_resultaten}

\subsection{Verwachte resultaten}
De verwachte uitkomst van deze bachelorproef is een technisch onderbouwde PoC voor Master Data Management (MDM) binnen Microsoft Fabric, die fungeert als objectieve benchmark tegenover de lopende Profisee-pilot bij IPCOM. Deze PoC levert een reproduceerbare pipeline op die (1) artikel- en leveranciersrecords uit Zwitserland, Oostenrijk en Noorwegen harmoniseert op basis van bestaande mappings, (2) records matcht via zowel deterministische regels als AI/fuzzy matching, (3) survivorship rules toepast om een golden record samen te stellen, en (4) de output publiceert voor consumptie en validatie in Power BI. 

Op inhoudelijk niveau wordt verwacht dat de PoC een aantoonbare reductie van duplicatie en inconsistentie oplevert. Dit wordt geoperationaliseerd via een vergelijking tussen een baseline-meting en de PoC-output, maar bovenal door een vergelijking met de resultaten van de Profisee-setup op dezelfde dataset. Daarnaast wordt inzichtelijk gemaakt hoe vaak en in welke mate een golden record de originele waarden van een specifiek bronsysteem wijzigt.

\subsection{Verwachte visualisaties en metingen}
Om de resultaten objectief te evalueren en de benchmark visueel te ondersteunen, worden in Power BI minstens de volgende grafieken uitgewerkt:

\begin{itemize}
  \item \textbf{Benchmark succesratio (Fabric vs. Profisee)}: Een vergelijkende visualisatie die toont hoeveel unieke golden records (en correcte matches) beide systemen genereren op basis van dezelfde inputdata.
  \item \textbf{Impact op bronsystemen}: Een visualisatie die per bronsysteem (Zwitserland, Oostenrijk, Noorwegen) het percentage toont waarbij de uiteindelijke golden record-waarde afwijkt van de originele, lokale data. Dit beantwoordt direct de nieuwe onderzoeksvraag.
  \item \textbf{Duplicatiegraad per systeem (voor/na)}: Staafdiagram met op de x-as de bronsystemen en op de y-as het percentage records dat als duplicaat wordt geïdentificeerd en geconsolideerd.
  \item \textbf{Volledigheid per attribuut (voor/na)}: Staafdiagram met op de x-as kritieke attributen (bv. EAN, artikelnaam) en op de y-as het percentage ingevulde waarden. Verwacht wordt dat dit percentage stijgt doordat waarden over bronnen heen worden gecombineerd.
  \item \textbf{Traceerbaarheidsoverzicht}: Tabelvisualisatie waarin per golden record wordt getoond welke bronrecords zijn gematcht, inclusief referentie naar de lokale ID's, zodat bronsystemen deze "foreign keys" kunnen overnemen.
\end{itemize}

\subsection{Meerwaarde voor de doelgroep}
De doelgroep (met name de data engineers binnen IPCOM) verkrijgt met deze resultaten drie concrete voordelen. Ten eerste biedt de PoC een objectieve "Build vs. Buy" business case: het toont aan of een in-house oplossing in Fabric kan concurreren met de licentiekosten en complexiteit van Profisee. Ten tweede fungeren de doorontwikkelde mappings, rule-sets en analyses (zoals de AI-matching op artikelnaam) als direct bruikbare input, ongeacht welk platform uiteindelijk in productie gaat. Ten derde profiteren alle aangesloten entiteiten van de gerealiseerde golden records doordat lokale systemen eenduidig kunnen verwijzen naar verrijkte master data.

\subsection{Verwachte conclusie}
Op basis van de probleemsituatie bij IPCOM wordt verwacht dat de Fabric-PoC een aanzienlijke, meetbare verbetering oplevert in datakwaliteit binnen de multi-ERP-context. Specifiek wordt verwacht dat de toevoeging van AI of fuzzy matching op artikelnamen een significante meerwaarde biedt bovenop louter deterministische regels (zoals EAN-matches). De uiteindelijke conclusie zal uitwijzen of de in-house ontwikkelde logica in Microsoft Fabric kwalitatief evenwaardig is aan de commerciële Profisee-pilot, en welke architectuur de meest schaalbare en kostenefficiënte route biedt voor de verdere uitrol binnen de IPCOM-groep.